\documentclass[12pt]{amsart} 
\usepackage{amsfonts}
\usepackage{amscd}
\usepackage{amsmath}
\usepackage{amssymb}
\usepackage{amsthm}
\usepackage{enumitem}
\usepackage{pgf,tikz,tikz-cd}
\usepackage{mathrsfs}
\usepackage[margin=1in]{geometry}
\usepackage{array}
\usepackage{lipsum}
\usepackage{dsfont}
\usepackage{stmaryrd}

%\graphicspath{ {./diagrams/} }

\input{./latex-preamble/cm-preamble.tex}

\newtheorem*{lemma}{Lemma}


\begin{document}

{Hartshorne} \hfill {2-1-2025}

\bigskip\bigskip

\begin{center}

{\sc Chapter II.2, Week 2}

\end{center}

\begin{center}

  {Noah Parker} %\footnote{Collaborated with Kalev Martinson and Frank :).}}
    
\end{center}

\bigskip

\begin{enumerate}
	\item[2.1.] {
      Let $A$ be a ring, $X=\Spec A$ its spectrum.
      Fix an element $f\in A$, $D(f)=V(f)^c\subset X$ the associated distinguished open set.
      Show that $(D(f),\O_X|_{D(f)})\cong \Spec A_f$ as locally ringed spaces.
    }
\end{enumerate}

\begin{proof}
  Let $Y=\Spec A_f$, $S=\{f^n\}_{n\geqs 1}$, and $\phi:A\rightarrow S^{-1}A=A_f$ be the natural map $a\mapsto a/1$.
  Note that $f^n\in \p$ prime for some $n\geqs 1$ iff $f\in \p$---that is, $S\cap \p=\0$ iff $\p\notin D(f)$.
  Thus, AM 3.11.iv gives us a bijection
  \begin{align*}
    D(f)~&\longleftrightarrow Y \\
    \p ~&\longrightarrow S^{-1}\p, \\
    \phi^*(\q)~&\longleftarrow \q.
  \end{align*}
  Clearly the map $\q\mapsto \phi^*(\q)$ is continuous, so it suffices to show that $\p\mapsto S^{-1}\p$ is continuous.

  Let $V(\b)\subset A_f$ be closed, and denote $\psi:\p\mapsto S^{-1}\p$. 
  AM 3.11.i gives us that $\b=S^{-1}\a$ is extended.
  Then $\p\supset \a$ implies $S^{-1}\p\supset S^{-1}\a=\b$ whenever $f\notin \a,\p$, so that $V(\a)\cap D(f)\subset \psi^{-1}(\b)$.
  Similarly, $S^{-1}\p\supset S^{-1}\a$ implies $\phi^*(S^{-1}\p)\supset \phi^*(S^{-1}\a)\supset \a$.
  % That is, for all $a\in \a$, there is some $x\in \p$ and $n\geqs 1$ such that $f^n(a-x)=0$.
  % Hence, $f^na\in \p$, so that $\p$ prime and $f\notin \p$ gives us $a\in \p$.%\footnote{As a sanity check, $f\notin \p$ also implies $f\notin \a$, so that $V(\a)\neq 0$!}
  We conclude that $\psi^{-1}(V(\b))=V(\a)\cap D(f)\subset D(f)$ is closed, so that $\psi=(\psi^*)^{-1}$ is continuous.

  Note that $(\O_X|_{D(f)})_\p\cong A_\p\cong (A_\p)_f\cong (A_f)_{\psi(\p)}\cong (\O_Y)_{\psi(\p)}$ for each $\p\not\ni f$.
  Call this isomorphism $\alpha_\p:(\O_X|_{D(f)})_\p\to (\O_Y)_{\psi(\p)}$, and let 
  \[
    \alpha:\sqcup_{\p\in D(f)}(\O_X|_{D(f)})_\p
    \to \sqcup_{\q\in Y}(\O_Y)_\q
  \]
  be defined by $\alpha|_{(\O_X|_{D(f)})_\p}=\alpha_\p$.
  Then we have a map
  \begin{align*}
    \phi^\sharp(U):\O_X|_{D(f)}(U)~&\to (\phi^*)_*\O_{Y}(U) \\
                        s~&\mapsto \alpha\circ s\circ\phi^*.
  \end{align*}
  We have that 
  \[
    \alpha\circ s|_V\circ \phi^*=\alpha\circ s\circ \phi^*|_{(\phi^*)^{-1}(V)}
  \]
  as functions on $(\phi^*)^{-1}(V)$ whenever $V\subset U\subset D(f)$ are open and $s\in \O_X(U)$, so $\phi^\sharp(U)$ gives a sheaf homomorphism $\phi^\sharp:\O_X|_{D(f)}\to \O_Y$.
  $\phi^\sharp_\p=\alpha_\p$ is an isomorphism of rings for all $\p\in D(f)$, so $\phi^\sharp$ is itself a sheaf isomorphism.
  Then $(f,f^\sharp)$ is an isomorphism of locally ringed spaces, since ring isomorphisms preserve maximality.
\end{proof}

\begin{enumerate}
	\item[2.2.] {
      Let $(X,\O_X)$ be a scheme, and let $U\subset X$ be any open subset.
      Show that $(U,\O_X|_U)$ is a scheme.
      We call this the \emph{induced scheme structure} on the open set $U$, and we refer to $(U,\O_X|_U)$ as the open subscheme of $U$.
    }
\end{enumerate}

\begin{proof}
  Choose a basic cover $\{D(f_i)\}$ of $U$. 
  Then each 
  \[
    (D(f_i),(\O_X|_U)_{D(f_i)})= (D(f_i),\O_X|_{D(f_i)})\cong\Spec A_{f_i}
  \]
  by 2.1. 
  We conclude that $\{D(f_i)\}$ is an affine open cover of $(U,\O_X|_U)$, so that the latter locally ringed space is a sheaf.
\end{proof}


\begin{enumerate}
	\item[2.3.] {
      $(*)$ A scheme $(X,\O_X)$ is \emph{reduced} if for every open set $U\subset X$, the ring $\O_X(U)$ has no nilpotent elements.
      \begin{enumerate}
        \item Show that $(X,\O_X)$ is reduced if and only if for every $p\in X$, the local ring $\O_{X,p}$ as no nilpotent elements.
        \item Let $(X,\O_X)$ be a scheme.
          Let $(\O_X)_{\red}$ be the sheaf associated to the presheaf $U\mapsto \O_X(U)_{\red}$, where for any ring $A$, we denote by $A_{\red}$ the quotient of $A$ by its nilradical $\fR(A)$. 
          Show that $(X,(\O_X)_{\red})$ is a scheme.
          We call it the \emph{reduced scheme} associated to $X$, and denote it by $X_{\red}$.
          Show that there is a morphism of schemes $X_{\red}\to X$, which is a homeomorphism on the underlying topological spaces.
        \item Let $f:X\to Y$ be a morphism of schemes, and assume that $X$ is reduced.
          Show that there is a unique morphism $g:X\to Y_{\red}$ such that $f$ is obtained by composing $g$ with the natural map $Y_{\red}\to Y$.
      \end{enumerate}
    }
\end{enumerate}

\begin{enumerate}
	\item[2.4.] {
      Let $A$ be a ring, and let $(X,\O_X)$ be a scheme.
      Given a morphism $f:X\to\Spec A$, we have an associated map $f^\sharp:\O_{\Spec A}\to f_*\O_X$ on sheaves.
      Taking Global sections, we obtain a homomorphism $A\to\Gamma(X,\O_X)$.
      Thus, there is a natural map
      \begin{align*}
        \alpha:\Hom_{\Sch}(X,\Spec A)\to \Hom_{\Ring}(A,\Gamma(X,\O_X)).
      \end{align*}
      Show that $\alpha$ is bijective.
    }
\end{enumerate}

\begin{proof}
  We begin by showing the result for $X=\Spec B$ affine.
  In this case, $\alpha$ has an inverse map given by 
  \begin{align*}
    \beta:\Hom_{\Ring}(A,B)&~\longrightarrow \Hom_{\Sch}(\Spec B,\Spec A) \\
    \vphi&~\longmapsto (f,f^\sharp),
  \end{align*}
  where 
  \begin{align*}
    f(\p)&=\vphi^{-1}(\p), \\
    f^\sharp(s)(\p)&=\vphi_{f(\p)}(s(f(\p))).
  \end{align*}
  It suffices to show that $\beta$ and $\alpha$ indeed compose to the identity.

  Fix $(f,f^\sharp)\in \Hom_{\Sch}(\Spec B,\Spec A)$, and denote 
  \[
    (\beta\circ\alpha)(f,f^\sharp) = (g,g^\sharp).
  \]
  Let $\phi_{f(\p)}:A\to A_{f(\p)}$ and $\phi_\p':B\to B_\p$ be the natural maps.
  The commutative square
  \[
    \begin{tikzcd}[column sep = large]
      A\ar[d,"\phi_{f(\p)}"']
      \ar[r,"f^\sharp(\Spec A)"]&B\ar[d,"\phi'_\p"] \\
      A_{f(\p)}\ar[r,"f^\sharp_{f(\p)}"']&B_\p
    \end{tikzcd}
  \] 
  gives us that 
  \begin{align*}
    g(\p)&=(f^\sharp(\Spec A))^{-1}(\p) \\
         &=(\phi'_\p\circ f^\sharp(X))^{-1}(m_\p) \\
         &=(f_{f(\p)}^\sharp\circ \phi_{f(\p)})^{-1}(m_\p) \\
         &= {\phi}_{f(\p)}^{-1}(m_{f(\p)}) \\
         &= f(\p),
  \end{align*}
  since $f^\sharp$ is a local homomorphism.
  Then
  \begin{align*}
    g^\sharp(s)(\p) &= f^\sharp(X)_{g(\p)}(s(g(\p))) \\
                    &= f^\sharp_{f(\p)}(s(f(\p))) \\ 
                    &=f^\sharp(s)(\p).
  \end{align*}
  implies $(g,g^\sharp)=(f,f^\sharp)$.

  Fix $\vphi:A\to B$, and denote $\beta(\vphi)=(f,f^\sharp)$, $(\alpha\circ \beta)(\vphi)= \psi$.
  Recall that $a\in A$ is identified with $\p\mapsto a_\p\in A_\p$, and similarly for $b\in B$.
  Thus, $\psi(a)=\vphi(a)$ iff they're equal as maps $\Spec B\to \sqcup_{\p\in \Spec B}B_\p$.
  The commutative square
  \[
    \begin{tikzcd}[column sep = large]
      A\ar[d]
      \ar[r,"\vphi"]&B\ar[d] \\
      A_{f(\p)}\ar[r,"\vphi_{f(\p)}"']&B_\p
    \end{tikzcd}
  \] 
  gives us
  \[
    \psi(a)(\p) = \vphi_{f(\p)}(a(f(\p))) = \vphi_{f(\p)}(a_{f(\p)}) = \vphi(a)(\p).
  \]
  Hence, $(\alpha\circ\beta)(\vphi)=\vphi$.
  We conclude that $\beta=\alpha^{-1}$, so $\alpha$ is a bijection.

  Now, we consider the general case of $X$ a scheme. 
  Define a map 
  \begin{align*}
    \beta:\Hom_{\Ring}(A,\Gamma(X,\O_X))&~\longrightarrow \Hom_{\Sch}(X,\Spec A) \\
    \vphi&~\longmapsto (f,f^\sharp)
  \end{align*}
  as follows.
  Fix an affine cover $\{U_i=\Spec B_i\}$ of $X$, and define $f:X\to \Spec A$ by $f(x)=\vphi^{-1}(\p_x)$, where $\p_x$ is the inverse image in $\Gamma(X,\O_X)$ of the unique maximal ideal $m_x\trianglelefteq(\O_X)_x$.
  Naturality of the restriction maps gives us that $f|_{U_i}=f_i$, where $(f_i,f_i^\sharp)=\beta(\rho_{U_i}\circ \vphi)$ is defined in the affine case.
  
  Next, we define $f^\sharp$ as follows.
  Fix $V\subset X$ open and $s\in \O_{\Spec A}(V)$.
  Consider the sections $f^\sharp_i(V\cap U_i)(s)\in f_*\O_X(V\cap U_i)$, with $f^\sharp_i$ as above.
  We have that 
  \begin{align*}
    f^\sharp_i(V\cap U_i\cap U_j)(s)(\p) 
    = (\rho_{V\cap U_i\cap U_j}\circ \vphi)_{f(\p)} (s(f(\p)))
    = f^\sharp_j(V\cap U_i\cap U_j)(s)(\p) 
  \end{align*}
  for all $\p\in f^{-1}(V\cap U_i\cap U_j)$, so the sections agree on their overlap.
  Then they glue to a unique section $f^\sharp(V)(s)\in f_*\O_X(V)$.
  
  It suffices to verify that $\alpha$ and $\beta$ compose to the identity.
  Fix $(f,f^\sharp)\in \Hom_{\Sch}(X,\Spec A)$, and denote $(\beta\circ \alpha)(f,f^\sharp)=(g,g^\sharp)$.
  Then $g|_{U_i} = g_i$, where 
  \[
    (g_i,g_i^\sharp)= \beta(\rho_{U_i}\circ f^\sharp(\Spec A)) = \beta(f^\sharp(U_i)).
  \]
  This gives us $g_i=f|_{U_i}$ by the affine case, so that $g=f$. 

  Fix $V\subset \Spec A$ open and $s\in \O_{\Spec A}(V)$. 
  Note that the $g^\sharp(V)(s)$ is the gluing of the sections $g^\sharp(V\cap U_i)(s)$ defined by
  \begin{align*}
    g^\sharp(V\cap U_i)(s)(\p)
    &=(\rho_{V\cap U_i}\circ f^\sharp(V))_{g(\p)} (s(g(\p))) \\
    &= f^\sharp_{f(\p)}(s(f(\p))) \\
    &= f^\sharp(V\cap U_i)(s)(\p)
  \end{align*}
  for all $\p\in f^{-1}(V\cap U_i)$.
  Thus, $(f,f^\sharp)=(g,g^\sharp)$.

  Finally, fix $\vphi:A\to \Gamma(X,\O_X)$ and let $(f,f^\sharp)= \beta(\vphi)$, $\psi=(\alpha)(f,f^\sharp)$.
  Then 
  \[
    \rho_{U_i}\circ \psi = \rho_{U_i}\circ (\alpha(f,f^\sharp)) = f^\sharp(U_i)= \rho_{U_i}\circ \vphi
  \]
  by the affine case, so that $\psi = \vphi$ by uniqueness of gluing.
  We conclude that $\alpha$ is a bijection.
  % We have that the following diagram commutes.
  % \[
  %   \begin{tikzcd}[row sep = 0.25cm]
  %     &&\Gamma(U_i,\O_X)\ar[dr]\\
  %     A\ar[r,"\vphi"]& \Gamma(X,\O_X)\ar[dr]\ar[ur] &&\Gamma(U_i\cap U_j,\O_X)\\
  %                    &&\Gamma(U_j,\O_X)\ar[ur]
  %   \end{tikzcd}
  % \] 
  % Noting that $\p_x\in \Spec B_x$ is the inverse image of $\p_x
  % Fix an affine cover $\{U_i\cong \Spec B_i\}$ of $X$.
  % Suppose $f:X\to \Spec A$ is a morphism of schemes.
  % Then it restricts to a set of morphisms $f_i:\Spec B_i\to \Spec A$ of affine schemes.
  % If $f,g:X\to \Spec A$ are two morphisms of schemes with $\alpha(f)=\alpha(g)$, 
\end{proof}

\begin{enumerate}
	\item[2.5.] {
      Describe $\Spec \ZZ$, and show that it is a final object for the category of schemes, i.e., each scheme $X$ admits a unique morphism to $\Spec\ZZ$.
    }
\end{enumerate}

\begin{proof}
  The only primes of $\ZZ$ are $(p)$ for $p$ prime, and $(0)$.
  The former are all maximal, while the latter is a generic point for $\Spec \ZZ$, since $V(\a)\ni (0)$ implies $(0)\supset \a$, i.e. $\a = (0)\subset \p$ for all $\p$ prime.

  We have a bijection
  \begin{align*}
    \Hom_{\Sch}(X,\Spec \ZZ)\to \Hom_{\Ring}(\ZZ,\Gamma(X,\O_X)),
  \end{align*}
  so that the former is singleton for all $X$ by the fact that $\ZZ$ is initial in $\Ring$.
\end{proof}

\begin{enumerate}
	\item[2.7.] {
      Let $X$ be a scheme.
      For any point $x\in X$, let $\O_x$ be the local ring at $x$, and $m_x$ its maximal ideal.
      We define the \emph{residue field} of $x$ on $X$ to be the field $k(x)=\O_x/m_x$.
      Now, let $K$ be any field.
      Show that to give a morphism of $\Spec K$ to $X$, it is equivalent to giving a point $x\in X$ and an inclusion map $k(x)\to K$.
    }
\end{enumerate}

\begin{proof}
  Note that $\Spec K=\{(0)\}$ is singleton, so a morphism $f:\Spec K\to X$ is a point $x=f((0))\in X$, as well as the information $f^\sharp(U): \O_X(U)\to K$ for every $x\in U\subset X$ open.
  Restriction is the identity on $\O_{\Spec K}$, so this defines a map $f^\sharp_x:\O_x\to K$.
  Since each element of $m_x$ is not invertible, this gives us a map $\ol{f^\sharp_x}:k(x)\to K$.
  It is a morphism of fields, and thus an injection.

  Conversely, suppose we have an inclusion $k(x)\hookrightarrow K$.
  Then, for each $x\in U\subset X$, we have a map 
  \[
    f^\sharp(U):\O(U)\to \O_x\to k(x)\hookrightarrow K.
  \]
  These maps play well with restriction, thus giving us a morphism $f^\sharp:\O_X\to \O_{\Spec K}$ of sheaves.
  Defining $f:\Spec K\to X:(0)\mapsto x$ gives us a morphism $(f,f^\sharp):\Spec K\to X$ of schemes, as desired.
\end{proof}

\begin{enumerate}
	\item[2.8.] {
      $(*)$ Let $X$ be a scheme.
      For any point $x\in X$, we define the \emph{Zariski tangent space $T_x$} at $x\in X$ to be the dual of the $k(x)$-vector space $m_x/m_x^2$.
      Now, assume that $X$ is a scheme over a field $k$, and let $k[\epsilon]/\epsilon^2$ be the \emph{ring of dual numbers} over $k$.
      Show that to give a $k$-morphism of $\Spec k[\epsilon]/\epsilon^2$ to $X$ is equivalent to giving a point $x\in X$, \emph{rational over $k$} (i.e. such that $k(x)=k$), and an element of $T_x$.
    }
\end{enumerate}

\begin{enumerate}
  \item[2.9.] {
      If $X$ is a topological space, and $Z$ an irreducible closed subset of $X$, a generic point for $Z$ is a point $\zeta$ such that $Z=\{\zeta\}^-$.
      If $X$ is a scheme, show that every (nonempty) irreducible closed subset has a unique generic point.
    }
\end{enumerate}

\begin{proof}
  We begin by proving the case of $X=\Spec A$ affine.
  Let $Z=V(\q)\subset X$ be irreducible and closed.
  That is, $V(\q)=V(\a)\cup V(\b)=V(\a\b)$ implies $V(\a)=V(\q)$ or $V(\b)=V(\q)$.
  Suppose now that $ab\in \q$.
  Then $V((\q,a)(\q,b))= V(\q)$ gives us $(\q,a)\subset \sqrt{(\q,a)} =\q$ or $(\q,b)\subset \sqrt{(\q,b)}=\q$.
  Hence, either $a\in \q$ or $b\in \q$, so that $\q$ is prime.

  % Conversely, take $\q\trianglelefteq A$ prime, and suppose that $V(\q)=V(\a)\cup V(\b)$, so $\q=\sqrt{\a}\sqrt{\b}=\sqrt{\a},\sqrt{\b}$.
  % Then $\q$ prime gives us $\q\supset \sqrt{\a}$ or $\q\supset \sqrt{\b}$, as otherwise there exists $ab\in \q$ with $a,b\notin \q$.
  % Thus, $V(\q)$ is irreducible.

  Suppose that $\zeta$ is a generic point associated to $Z=V(\q)$.
  Then 
  \begin{align*}
    \{\zeta\}^- &= \bigcap_{V(\a)\ni \zeta}V(\a) 
                = V(\cup_{\a\subset \zeta}\a) 
                = V(\zeta)
  \end{align*}
  implies that $\zeta = \q$. 
  This is clearly unique, and we showed that $\q$ is prime above, so this is always possible.

  Let $Z$ be an irreducible closed subset of a scheme $X$, and let $\{U_i\cong \Spec A_i\}_{i\in I}$ be an affine open cover of $X$.
  Replace $\{U_i\}$ by a subcover such that each $U_i\cap Z\neq \0$.
  We claim that each $U_i\cap Z$ is irreducible in $U_i$.
  Take $\0\neq U_i\cap V\subset U_i\cap Z$ nonempty open in $U_i$.
  Then $\0\neq U_i\cap V\subset Z$ nonempty open gives us that 
  \[
    \Cl_{U_i}(U_i\cap V) = U_i\cap \Cl_X(U_i\cap V)= U_i\cap Z,
  \]
  whence $U_i\cap Z$ is irreducible in $U_i$.

  Take $\zeta_i\in U_i$ to be the unique generic point of $U_i\cap Z$ given by the affine case, so $\{\zeta_i\}^-\supset \Cl_{U_i}(\zeta_i)=Z\cap U_i$.
  Note that $Z\cap U_i\cap U_j\neq 0$ for all $i,j$, as otherwise 
  \[
    Z= Z\cap (U_i^c\cup U_j^c)= (Z\cap U_i^c)\cup(Z\cap U_j^c)
  \]
  gives us that $Z$ is disjoint from $U_i$ or $U_j$.
  Then 
  \[
    \{\zeta_i\}^-\cap U_j \supset (Z\cap U_i)\cap U_j\neq \0
  \] 
  implies $\zeta_i\in U_j$, as otherwise $U_j^c\cap \{\zeta_i\}^-\supset \{\zeta_i\}^-$.
  Then $\Cl_{U_j}(\zeta_i)=\{\zeta_i\}^-\cap U_j$ implies 
  \[
    \Cl_{U_j}(\zeta_i)\supset (Z\cap U_i\cap U_j)^-\cap U_j = Z\cap U_j.
  \]
  Hence, $\zeta_i$ is a generic point of $Z\cap U_j\subset U_j$ for each $j$.
  Uniqueness in the affine case gives us $\zeta_i=\zeta_j=:\zeta$ for all $i,j$, so that
  \[
    \{\zeta\}^-\supset \bigcup_{i\in I} U_i\cap Z = Z.
  \] 
  Any generic point $\zeta'$ of $Z$ must also be a generic point of each of the $U_i\cap Z$, so uniqueness follows from the affine case.
\end{proof}

\begin{enumerate}
  \item[2.10.] {
      Describe $\Spec \R[x]$. How does its topological space compare to the set $\R$? To $\CC$?
    }
\end{enumerate}

\begin{proof}
\end{proof}

\begin{enumerate}
  \item[2.13.] {
      $(*)$ A topological space is \emph{quasi-compact} if every open cover has a finite sub-cover.
      \begin{enumerate}
        \item Show that a topological space $X$ is noetherian if and only if every open subset is quasi-compact.
        \item If $A$ is a noetherian ring, show that $\Sp(\Spec A)$ is a noetherian topological space.
        \item Give an example to show that $\Sp(\Spec A)$ can be noetherian even when $A$ is not.
      \end{enumerate}
    }
\end{enumerate}

\begin{proof}
\end{proof}

\begin{enumerate}
	\item[2.14.] {
      \begin{enumerate}
        \item Let $S$ be a graded ring.
          Show that $\Proj S=\0$ if and only if every element of $S_+$ is nilpotent.
        \item Let $\vphi:S\to T$ be a graded homomorphism of graded rings (preserving degrees).
          Let $U=\{\p\in \Proj T~:~\p\not\supset\vphi(S_+)\}$.
          Show that $U$ is an open subset of $\Proj T$, and show that $\vphi$ determines a natural morphism $f:U\to \Proj S$.
        \item The morphism $f$ can be an isomorphism even when $\vphi$ is not.
          For example, suppose that $\vphi:S_d\to T_d$ is an isomorphism for all $d\geqs d_0$, where $d_0$ is any integer.
          Then show that $U=\Proj T$ and the morphism $f:\Proj T\to \Proj S$ is an isomorphism.
      \end{enumerate}
    }
\end{enumerate}

\begin{proof}
  (a) Denote $\nil' S := \cap_\p\p$, where $\p$ ranges over all homogeneous prime ideals of $S$.
  It suffices to show that $\nil' S=\nil S$.
  % First, we show that homogeneous prime ideals are prime in the usual sense.

  % Fix a homogeneous prime ideal $\p$.
  % Suppose that $ab\in \p$, and take $a=\sum_{i=1}^ra_i$, $b=\sum_{j=1}^sb_j$.
  % Then 
  % \[
  %   ab = \sum_{\substack{1\leqs i\leqs r \\ 1\leqs j\leqs s}}a_ib_j\in\p
  % \] 
  % gives us $a_rb_s\in \p$, so that $a_r\in\p$ or $b_s\in \p$.
  % Without loss of generality, $a_r\in \p$, implying $(a-a_r)b = ab-a_rb\in \p$.
  % Induction allows us to reduce to the $a,b$ homogeneous of degree 0 case, whence $\p$ is prime.
  % It follows that $\nil S\subset \nil' S$.

  % Suppose $f\in \nil S$.
  % Write $f=\sum_{i=1}^d f_i$, and take $n>0$ with $f^n=0$.
  % Then $f_d^n=0$ gives us that $f_d\in \nil S$.
  % Replacing $f$ by $f-f_d\in \nil S$ and repeating the process gives us that each $f_i\in \nil S$.
  % Thus, $\nil S$ is a homogeneous prime ideal.
  For any prime ideal $\p$, the homogenization $h(\p)\subset \p$ generated by its homogeneous elements is clearly a homogeneous prime ideal.
  Then we have
  \[
    \nil' S \supset \nil S=\bigcap_{\p}\p \supset \bigcap_{\p}h(\p)=\nil' S,
  \] 
  %since each homogeneous prime ideal is prime in the usual sense.
  which gives us $\nil'S\subset \nil S$.
  We conclude that $\nil S=\nil' S$.
  Now, $\Proj S=\0$ iff $S_+\subset \nil' S=\nil S$, which is exactly our result.

  (b) Let $\p\in \Proj T$ be a homogeneous prime ideal.
  Then $\vphi^{-1}(\p)$ is a homogeneous prime ideal, and $\vphi^{-1}(\p)\supset S_+$ implies 
  \[
    \p\supset \vphi(\vphi^{-1}(\p))\supset \vphi (S_+),
  \]
  i.e. $\p\notin U$.
  Conversely, $\p\supset \vphi(S_+)$ implies 
  \[
    \vphi^{-1}(\p)\supset \vphi^{-1}(\vphi(S_+))\supset S_+,
  \]
  so that $\vphi^{-1}(\p)\notin \Proj S$.
  We conclude that $U$ is the inverse image of $\Proj S$ under the mapping of homogeneous prime ideals $\p\mapsto \vphi^{-1}$.
  Hence, we have a natural map $f:U\to\Proj S:\p\mapsto \vphi^{-1}(\p)$.

  Fix a homogeneous $g\in T_+$.
  Then
  \begin{align*}
    \vphi^{-1}(\p)\ni g\implies \p\supset \vphi(\vphi^{-1}(\p))\ni \vphi(g) \implies \vphi^{-1}(\p)\supset \vphi^{-1}(\vphi(g))\ni g
  \end{align*}
  gives us that
  \begin{align*}
    f^{-1}(D_+(g)) &= D_+(\vphi(g))\subset U.
  \end{align*}
  Hence, $f:U\to \Proj S$ is continuous.
  We must now define the map $f^\sharp:\O_S\to f_*(\O_T|_U)$ of sheaves.

  It suffices to define $f^\sharp$ on basic open sets.
  Fix a homogeneous $g\in S_+$, and note that 
  \[
    f_*(\O_T|_U)(D_+(g)) = \O_T(f^{-1}(D_+(g))\cap U) = \O_T(D_+(\vphi(g)))\cong T_{(\vphi(g))}.
  \]
  Recall as that $\vphi$ is degree preserving, so each of its localizations are as well.
  In particular, $\vphi_g:S_g\to T_g$ restricts to a map $\vphi_g:S_{(g)}\to T_{(\vphi(g))}$.
  Then the following commutative diagram gives us our sheaf map on $D_+(g)$.
  \[
    \begin{tikzcd}[column sep = 1.2cm]
      \O_S(D_+(g))\ar[d, "\sim" {rotate=90, anchor=north}]\ar[r,"f^\sharp(D_+(g))"]
      &f_*(\O_T|_U)(D_+(g))\ar[d, "\sim" {rotate=90, anchor=north}] \\
      S_{(g)}\ar[r,"\vphi_g"']&T_{(\vphi(g))}
    \end{tikzcd}
  \] 
  To show that these maps glue together, we note that $D(g)\cap D(h)= D(gh)$, and consider the following commutative diagram.
  \[
    \begin{tikzcd}[column sep = 0.05cm,row sep = 0.4cm]
      S_{(g)}\ar[rr,"\vphi_g"]\ar[dr]
        &&S_{(\vphi(g))}\ar[dr]\\
      &S_{(gh)}\ar[rr,"\vphi_{gh}"] 
        &&S_{(\vphi(gh))} \\
      S_{(h)}\ar[ur]\ar[rr,"\vphi_h"']
        &&S_{(\vphi(h))}\ar[ur]
    \end{tikzcd}
  \] 
  This defines our sheaf homomorphism on all of $\Proj S$, so that we're done.

  % Let $s\in \O_S(V)$, and define 
  % \[
  %   f^\sharp(V)(s)(\p) = \vphi_{f(\p)}(s(f(\p)))
  % \]
  % for all $\p\in f^{-1}(U\cap V)\subset \Proj T$.
  % $\vphi$ is degree preserving, so each of its homogeneous localizations $\vphi_{f(\p)}$ is degree preserving.

  % If $s|_W(\q)=a/g\in (\O_S)_{(\q)}$ for $a,g\in S$ homogeneous of the same degree and $g\notin \q$ for all $\q\in U\cap W$, then 
  % \[
  %   f^\sharp(W)(s|_W)(\p) = \vphi_{f(\p)}(a/g)\in \vphi_{f(\p)}(\O_S)_{(f(\p))}\subset (\O_T)_{(\p)}
  % \]
  % for all $\p\in f^{-1}(U\cap W)$.
  % $\vphi_{f(\p)}(a/g)=\vphi(a)/\vphi(g)$, with $\vphi(a),\vphi(b)\in T$ homogeneous of the same degree since $\vphi$ is degree preserving, and $\vphi(g)\notin \vphi(f(\p))\subset \p$.
  % Each $f^\sharp_\p=\vphi_\p:(\O_S)_{(f(\p))}\to (\O_T)_{(\p)}$ is a local homomorphism, since $\vphi^{-1}(\p)=f(\p)$ by definition.
  % Naturality of $f^\sharp$ with respect to restriction is trivial, so we conclude that $(f,f^\sharp)$ is indeed a morphism of locally ringed spaces.

  (c) Fix $d_0$ such that $\vphi:S_d\to T_d$ is an isomorphism for all $d\geqs d_0$.
  We wish to show that $U=\Proj T$, so it suffices to show that $\sqrt{\vphi(S_+)}=T_+$.
  One inclusion is clear since $\vphi$ is degree preserving.
  Surjectivity of $\vphi:S_d\to T_d$ gives us $\vphi(S_+)\supset \oplus_{d\geqs d_0}T_d$.
  Then, for any element $a\in T_+$, $a^{d_0}$ is of degree at least $d_0$, so $a^{d_0}\in \vphi(S_+)$ implies $a\in \sqrt{\vphi(S_+)}$.
  We conclude that $\sqrt{\vphi(S_+)}=T_+$, so $U=\Proj T$.

  % It suffices to show that each map of stalks $\vphi_{f(\p)}:S_{(f(\p))}\to T_{(\p)}$ is an isomorphism.
  % Recall that $\vphi_{f(\p)}(a/g)=\vphi(a)/\vphi(g)$, whenever $a,g\in S$ are homogeneous of the same degree and $g\notin f(\p)$.
  % $f(\p)\in \Proj S$ gives us $h\in S_+\setminus f(\p)$.
  % Replacing $a/g$ by $ah^{d_0}/gh^{d_0}$, we may assume that $a,g$ are of degree $\geqs d_0$.
  % Then
  % \begin{align*}
  %   \vphi_{f(\p)}(a/g) &= \vphi_{f(\p)}(a'/g')\\
  %   \implies ~\frac{\vphi(a)}{\vphi(g)}
  %                      &=\frac{\vphi(a')}{\vphi(g')}.
  % \end{align*}
  % Thus, we can find $t\in T_+\setminus \p$ such that
  % \[
  %   t(\vphi(a)\vphi(g')-\vphi(a')\vphi(g))=0.
  % \] 
  % Without loss of generality, we may assume $t\in T_d$ for some $d\geqs d_0$.
  % Then $T_d=\vphi(S_d)$ gives us $s\in S_d$ with $t=\vphi(s)$, so
  % \begin{align*}
  %   \vphi(t(ag'-a'g))=0\implies \frac{a}g=\frac{a'}{g'}
  % \end{align*}
  % by injectivity of $\vphi:\oplus_{d\geqs d_0}S_d\to \oplus_{d_\geqs d_0}T_d$.
  % Thus, $\vphi_{f(\p)}:S_{(f(\p))}\to T_{(\p)}$ is injective.

  % Take $b/h\in T_{(\p)}$. 
  % By the above argument, we may assume $b,h$ are of degree $d\geqs d_0$.
  % Then we can take homogeneous $a,g$ of degree $d\geqs d_0$ such that $\vphi(a)=b$ and $\vphi(g)=h$, so that
  % \begin{align*}
  %   \vphi_{f(\p)}(a/g)= \frac{\vphi(a)}{\vphi(g)} = \frac{b}h
  % \end{align*}
  % gives surjectivity of $\vphi_{f(\p)}$.
  % We conclude that $\vphi_{f(\p)}$ is an isomorphism for each $\p\in \Proj T$, so that $\vphi$ is an isomorphism of schemes.

  It suffices to show that $f^\sharp$ is an isomorphism on open sets $D(g)$.
  Fix homogeneous $g\in S_+$.
  Since $D(g^n)=D(g)$ for all $n>0$, we may assume $g$ is of degree $\geqs d_0$.
  We show the map $\vphi_g:S_{(g)}\to T_{(\vphi(g))}$ is surjective.
  For any $b/\vphi(g)^n\in (B_A)_\p$, $b\vphi(g)=\vphi(a)$ with $a\in A$ gives us 
  \[
    \vphi_g(a/g^{n+1})=\vphi(ag)/\vphi(g^{n+1})=b/\vphi(g)^n\in \im\vphi_g.
  \]
  It suffices to check injectivity.
  If $\vphi(a/g)=0$, then we can find $n\geqs 1$ such that $\vphi(g)^n\vphi(a)=0$.
  Then $\vphi(g^na)=0$ implies $g^na=0$ by injectivity of $\vphi$, whence $a/g =0$ in $S_{(g)}$.
  The result follows.
\end{proof}


\begin{enumerate}
  \item[2.16.] {
      Let $X$ be a scheme, let $f\in \Gamma(X,\O_X)$, and define 
      \[
        X_f:=\{x\in X: f_x\notin \m_x\subset \O_{X,x}\}.
      \]
      \begin{enumerate}
        \item If $U=\Spec B$ is an open affine subscheme of $X$, and if $\ol{f}\in B=\Gamma(U,\O_X|_U)$ is the restriction of $f$, show that $U\cap X_f=D(\ol{f})$.
          Conclude that $X_f$ is an open subset of $X$.
        \item Assume that $X$ is quasi-compact. Let $A=\Gamma(X,\O_X)$, and let $a\in A$ be an element whose restriction to $X_f$ is $0$.
          Show that, for some $n>0$, $f^na=0$.
        \item Now assume that $X$ has a finite cover by open affines $U_i$ such that each intersection $U_i\cap U_j$ is quasi-compact. 
          (This hypothesis is satisfied, for example, if $\ssp(X)$ is noetherian.)
          Let $b\in \Gamma(X_f,\O_{X_f})$.
          Show that, for some $n>0$, $f^nb$ is the restriction of an element of $A$.
        \item With the hypothesis of (c), conclude that $\Gamma(X_f,\O_{X_f})\cong A_f$.
      \end{enumerate}
    }
\end{enumerate}

\begin{proof}
  (a) $\p\in U\cap X_f$ iff $\ol{f}_\p\notin \m_\p=\p\O_{U,\p}$, i.e. $\ol{f}\notin \p$.
  It follows that $U\cap X_f= D(\ol{f})$ is open, so that $X_f$ is a union of open sets, hence open itself.

  (b) Let $a\in A$ with $a|_{X_f}=0$.
  Then, for any open affine $U$ intersection $X_f$, we define denote the restriction $\ol{a}:=a|_U$.
  Then $\ol{a}|_{D(\overline{f})}=a|_{U\cap X_f}=0$, whence there exists $n >0$ such that $\ol{f^na}=\ol{f}^n \ol{a} =0$.
  Quasi-compactness of $X$ gives us a finite affine open cover $\{U_i\}$, so that we can take some fixed $n>0$ with $(f^na)|_{U_i}=0$ for all $U_i$.
  The sheaf axiom gives us that in fact $f^na=0$, so we are done.

  (c) Consider $f_i:=f|_{U_i}$ and $b_i:=b|_{D(f_i)}$.
  Then, for each $i$, (a) gives us some $m>0$ such that 
  \[
    (f_i^mb_i)|_{D(f_i)}=c_i|_{D(f_i)},
  \]
  with $c_i\in \Gamma(U_i,\O_{U_i})$.
  Denote $U_{ij}=U_i\cap U_j$.
  $(c_i-c_j)|_{U_{ij}\cap X_f} = 0$, so (c) gives us $l>0$ such that $f^l(c_i-c_j)|_{U_ij} = 0$, i.e. $(f^lc_i)|_{U_{ij}}=(f^lc_j)|_{U_{ij}}$.
  There are only finitely many $U_{ij}$, so we can take $l$ such that the above holds for each $i,j$. 
  Then we have $n= m+l>0$ such that 
  \begin{align*}
    (f_i^nb_i)|_{D(f_i)} &= c_i|_{D(f_i)}, \\ 
    c_i|_{U_ij} & = c_j|_{U_{ij}}
  \end{align*}
  for all $i,j$.
  Then the $c_i$ glue to an element $c\in A$ such that $c|_{X_f}$ is the gluing of the $f_i^nb_i|_{D(f_i)}$; namely, $f^nb$.
  
  (d) By (c), we have a map $\phi:\Gamma(X_f,\O_{X_f})\to A_f:b\mapsto f^nb/f^n$, which is well defined regardless of our choice of $n>0$.\footnote{We are often somewhat sloppy in writing $f^nb$ for the element of $A$ which restricts to it, but we attempt to distinguish when necessary.}
  Suppose that $\phi(b)=0$, so that there exists $m>0$ with $f^{n+m}b =f^m(f^nb\cdot 1 - f^n\cdot 0) = 0$.
  Then, with notation as in (c), $f_i^{n+m}b_i=0$ means $b_i=0$ for each $i$, whence $b=0$, showing injectivity.

  Finally, we show surjectivity.
  Fix $a/f^m\in A_f$.
  It suffices to find $b\in \Gamma(X_f,\O_{X_f})$ such that $f^{n+m}b-f^na=0$ for $f^nb\in A$.\footnote{It is in fact equivalent, as if $f^{n+m+l}b-f^{n+l}a=0$ for $l>0$, then we can replace $n>0$ by $n+l>0$.}
  Fortunately, we need only glue together the maps $b_i = a|_{D(f_i)}/f_i^m$ to a map $b\in \Gamma(X_f,\O_{X_f})$, where compatibility is simple.
  Then $f^mb=a$, so that $\phi(b)=f^mb/f^m=a/f^n$, as desired.
  We conclude that $\phi$ is an isomorphism of rings.
\end{proof}

\begin{enumerate}
  \item[2.17.] {
      \emph{A Criterion for Affineness}
      \begin{enumerate}
        \item Let $f:X\to Y$ be a morphism of schemes, and suppose that $Y$ can be covered by open subsets $U_i$ such that for each $i$, the induced map $f^{-1}(U_i)\to U_i$ is an isomorphism.
          Then $f$ is an isomorphism.
        \item A scheme $X$ is affine if and only if there is a finite set of elements $f_1,\ldots, f_r\in A=\Gamma(X,\O_X)$ such that the open subsets $X_{f_i}$ are affine, and $f_1,\ldots, f_r$ generate the unit ideal in $A$.
      \end{enumerate}
    }
\end{enumerate}

\begin{proof}
  (a) Let 

  (b) If $X$ is affine, then we can take $(f_1,\ldots, f_r)=(1)$ to be a finite set of generators of the unit ideal, so that $X_{f_i}=D(f_i)$ are distinguished affines covering $X$.
  
  Conversely, take a scheme $X$ with such elements $f_1,\ldots, f_r\in A$.
  Each $X_{f_i}$ is affine and thus quasi-compact, so a finite cover of $X$ by quasi-compact open sets implies $X$ is quasi-compact.
  Each intersection $X_{f_i}\cap X_{f_j}=D_{X_{f_i}}(f_j)\cong \Spec (A_{f_i})_{f_j}\cong \Spec A_{f_if_j}$ is affine with and thus quasi-compact, so that 2.16(d) gives us $\Gamma(X_{f_i},\O_{X_{f_i}})\cong A_{f_i}$.

  Ex. 2.4 supplies us with an isomorphism $\phi_i:X_{f_i}\iso\Spec A_{f_i}$, so that we have a cover of $\Spec A$ by $X_{f_i}\cong \Spec A_{f_i}$ where $\phi_i|_{D_{X_{f_i}(f_j)}}=\phi_j|_{D_{X_{f_j}}(f_i)}$.
  Then we can glue together the maps $\phi_i$ to a map $\phi:X\to \Spec A$, which is an isomorphism by (a).
  Hence, $X\cong \Spec A$ is affine.
\end{proof}


\begin{enumerate}
	\item[2.18.] {
      In this exercise, we compare some properties of a ring homomorphism to the induced morphism of the spectra of the rings.
      \begin{enumerate}
        \item Let $A$ be a ring, $X=\Spec A$, and $f\in A$.
          Show that $f$ is nilpotent if and only if $D(f)$ is empty.
        \item Let $\vphi:A\to B$ be a homomorphism of rings, and let $f:Y=\Spec B\to X=\Spec A$ be the induced morphism of affine schemes.
          Show that $\vphi$ is injective if and only if the map of sheaves $f^\sharp:\O_X\to f_*\O_Y$ is injective.
          Show furthermore in that case $f$ is \emph{dominant}, i.e., $f(Y)$ is dense in $X$.
        \item With the same notation, show that if $\vphi$ is surjective, then $f$ is a homeomorphism of $Y$ onto a closed subset of $X$, and $f^\sharp:\O_X\to f_*\O_Y$ is surjective.
        \item Prove the converse to (c).
          Namely, if $f:Y\to X$ is a homeomorphism onto a closed subset, and $f^\sharp:\O_X\to f_*\O_Y$ is surjective, then $\vphi$ is surjective.
      \end{enumerate}
    }
\end{enumerate}

\begin{proof}
  (a) The result is trivial, as $D(f)=\0$ iff $f\in \cap_\p\p=\nil f$.

  (b) In the language of coherent sheaves, we have $\O_X = \wtilde{A}$ and $f_*\O_Y=f_*\wtilde{B} = \wtilde{B_A}$.
  Then $f^\sharp_\p:\O_{X,\p}\to (f_*\O_Y)_\p$ corresponds to the map $\vphi_\p:A_\p\to (B_A)_\p$.
  By CA, we have that $\vphi$ is injective iff each $\vphi_\p$ is injective, which by the above occurs iff $f^\sharp:\O_X\to f_*\O_Y$ is injective.

  It suffices to show that $f(Y)$ is dense in $X$.
  Fix $g\in A$.
  Then
  \begin{align*}
    f(Y)\cap D(g)=\0 &\iff f^{-1}(D(g))=\0 \\
                     &\iff D(\vphi(g))=\0 \\
                     &\iff \vphi(g)\in \nil B \\
                     &\iff g\in \nil A \\
                     &\iff D(g)=\0
  \end{align*}
  by (a) and injectivity of $\vphi$ respectively.
  Any nonempty open $U\subset X$ contains a nonempty distinguished open of $X$, so $f(Y)$ is dense in $X$.
  
  % Clearly $f^\sharp$ injective implies $\vphi=f^\sharp(X)$ is injective.
  % Suppose conversely that $\vphi$ is injective.
  % It suffices to show that each map $f^\sharp_{f(\p)}=\vphi_{f(\p)}$ of stalks is injective, for $\p\in Y$.
  % 
  % % Let $a/g\in \ker\vphi_{f(\p)}$, so there is $t\in B\setminus \p$ with $t\vphi(a)=0$.
  % % If $t=\vphi(s)\in \vphi(A\setminus f(\p))\subset B\setminus \p$, then $0=t\vphi(a)=\vphi(sa)$ implies $sa=0$ in $A$, whence $a/g=0$ in $A_{f(\p)}$.
  % % Suppose instead that $\vphi(s)\vphi(a)\neq 0$ for all $s\in A\setminus f(\p)$.

  % We give $B$ the structure of an $A$-module by $a\cdot b = \vphi(a)b$.
  % Then $a\cdot \vphi(a')=\vphi(a)\vphi(a')=\vphi(aa')$ gives us that $\vphi$ is an injective morphism of $A$-module.
  % Flatness of the $A$-module $A_{f(\p)}$ gives us that $\vphi_{f(\p)}:A_{f(\p)}\hookrightarrow A_{f(\p)}\otimes_A B$ is injective.
  % We have an isomorphism $A_p\otimes_A B\xrightarrow{\sim}S^{-1}B$, with $S=\vphi(A\setminus f(\p))\subset B\setminus \p$.
  % Then we have an injection $S^{-1}B\hookrightarrow (S^{-1}B)_{S^{-1}\p}\cong B_\p$.
  % Composing these together, we see that
  % \[
  %   \vphi_{f(\p)}:
  %   A_{f(\p)}\hookrightarrow 
  %   A_{f(\p)}\otimes B \xrightarrow{\sim}
  %   S^{-1}B \hookrightarrow
  %   (S^{-1}B)_{S^{-1}\p}\xrightarrow{\sim}
  %   B_\p
  % \] 
  % is an injection.

  % It suffices to show that, in this case $f(Y)$ is dense in $X$.
  % Suppose $V(\a)\supset f(Y)$ is closed, and write $V(\a)^c = \cup_{g\in \a}D(g)$.
  % Fix $g\in \a$, and note that $D(g)\cap f(Y)=\0$.
  % Hence, $g\in f(\p)=\vphi^{-1}(\p)$ for all $\p\in X$.
  % Thus, $\vphi(g)\in\vphi(\vphi^{-1}(\p))\subset \p$ for all $\p\in X$ implies $\vphi(g)$ is nilpotent.
  % Taking $n>0$ with $\vphi(g)^n=0$, we see that $\vphi(g^n)=0$, whence $g^n=0$ by injectivity of $\vphi$.
  % We conclude that $g$ is nilpotent, so that $D(g)=\0$.
  % Hence, $V(\a)= X$, showing that $f(Y)^-=X$ is dense.

  (c) First, suppose that $\vphi:A\to B$ is surjective.
  Localization is exact, so each of the stalks $\vphi_\p:A_\p\to (B_A)_\p$ is surjective, whence $f^\sharp:\O_X\to f_*\O_Y$ is surjective.
  Surjectivity of $\vphi$ gives us
  \[
    \p\supset \a\implies\vphi^{-1}(\p)\supset \vphi^{-1}(\a)\implies \p=\vphi(\vphi^{-1}(\p))\supset \vphi(\vphi^{-1}(\a))=\a,
  \]
  so that 
  \begin{align*}
    f(V(\a)) &= \{\vphi^{-1}(\p)\in X ~:~ \p\supset \a\} \\
             &= f(Y)\cap V(\vphi^{-1}(\a)).
  \end{align*}
  In particular,
  \[
    f(Y)=f(Y)\cap V(\ker\vphi)\subset V(\ker\vphi).
  \] 

  Take $\q\in V(\ker\vphi)$, and suppose there exists $a\in \vphi^{-1}(\vphi(\q))\setminus \q$.
  Then there is some $a'\in \q$ with $\vphi(a)= \vphi(a')$, so that $a-a'\in \ker\vphi\setminus \q=\0$, a contradiction.
  Thus, $\vphi^{-1}(\vphi(\q))=\q$.

  We now show that $\vphi(\q)$ is prime.
  Suppose $bb'\in \vphi(\q)$, and take $a,a'\in A$ such that $\vphi(a)=b$ and $\vphi(a')=b'$.
  Then $aa'\in \vphi^{-1}(\vphi(\q))=\q$ implies $a\in \q$ or $a'\in \q$.
  Hence, $b\in \vphi(\q)$ or $b'\in \vphi(\q)$ gives us that $\vphi(\q)$ is prime, so that $g:V(\ker \vphi)\to Y:\q\to\vphi(\q)$ is a set-theoretic inverse to $f:Y\to V(\ker\vphi)\subset X$.
  In particular, $f(Y)=V(\ker\vphi)$.

  It suffices to show that $g:V(\ker\vphi)\to Y$ is continuous.
  We write
  \begin{align*}
    f(V(\a)) = f(V)\cap V(\vphi^{-1}(\a)) = V(\ker\vphi + \vphi^{-1}(\a)),
  \end{align*}
  so that $g^{-1}(V(\a))= V(\ker\vphi + \vphi^{-1}(\a))$ is closed in $X$.
  We conclude that $f$ is a homeomorphism onto its closed image $f(Y)=V(\ker\vphi)$.

  (d) Suppose now that $f:Y\to X$ is a homeomorphism onto a closed subset, and $f^\sharp:\O_X\to\O_Y$ is surjective.
  Let $Z=\Spec(A/\ker \vphi)$, so that the universal property of quotients supplies $\ol{\vphi}:A/\ker\vphi\to B$ such that $\vphi = \ol{\vphi}\circ \pi$.
  We denote the associated morphism of sheaves $\ol{f}:Y\to Z$.
  Note that $\pi:A\to A/\ker\vphi$ is surjective, so (c) gives us that the associated map $g:Z\to V(\ker\vphi)\subset X$ is a homeomorphism.

  $\ol{\vphi}$ is injective, so (b) implies $\ol{f}^\sharp:\O_Z\to \O_Y$ is injective, and $\ol{f}(Y)=f(Y)$ is dense.
  $f(Y)=f(Y)^-$ is closed, so $\ol{f}:Y\to Z$ is surjective.
  $f = g\circ \ol{f}$ with $f,g$ homeomorphisms onto their image implies $\ol{f}$ is a homeomorphism onto its image $\ol{f}(Y)=Z$.
  Additionally, $f^\sharp = \ol{f}^\sharp g^\sharp$ surjective implies $\ol{f}^\sharp$ is surjective, so that $\ol{f}^\sharp:\O_Z\to \O_Y$ is an isomorphism of sheaves.
  Then 
  \[
    \vphi =\ol{\vphi}\circ \pi = \ol{f}^\sharp(Z)\circ \pi
  \]
  is the composition of a surjective map with an isomorphism, and is therefore surjective itself.
\end{proof}


% \begin{enumerate}
%   \item[19.] {
%     }
% \end{enumerate}
% 
% \begin{proof}
% \end{proof}
% 
% \begin{enumerate}
%   \item[20.] {
%     }
% \end{enumerate}
% 
% \begin{proof}
% \end{proof}
% 
% \begin{enumerate}
%   \item[21.] {
%     }
% \end{enumerate}
% 
% \begin{proof}
% \end{proof}

\end{document}
